\clearpage
\section{Universal converse operational robustness}
\label{sec:copying-operational-converse-robustness}
This section closes the converse obligation left in
Section~\ref{sec:copying-operational-forward-robustness}. For both micro-beta
and reference creation, every eligible right-hand machine rewrite now has a
constructed matching left-hand rewrite and an admissible residual specimen.
The proof covers all nine original local rules, arbitrary stable operands,
recursive beta bodies and copyable payloads, and sharing trees with unused
branches crossing template fillings. No paired occurrence, matching rewrite,
residual-admissibility or machine-determinism hypothesis is added to the
universal theorem. The earlier proved machine laws are used as lemmas.
The terms of Proposition~12.2 and their Figure~16 translation are unchanged.

\termdef{operational-converse-robustness}
\begin{theorem}[Universal converse operational robustness]
Let $T$ be either original operational pre-template, and let $Q,P,U$ be cost
orders. The opposite of its binding-free specimen system satisfies robustness
with costs $(Q,P,U)$. Given an eligible specimen and a right-hand rewrite,
the proof constructs one left-hand rewrite and a pointed quasi-specimen in
the opposite system. It uses zero additional right-hand transitions and one
left-hand transition, so the local cost inequality is $Q(1+0,1)$.
\end{theorem}
Here eligibility is precisely the original robustness definition: both
specimen states are rooted, the focus is in rewrite mode and is not entering
a template hole. Binding freedom and closed expanded graphs come from the
original specimen type. These are the intended contextual restrictions,
not new residual certificates supplied by an application. The proof uses
the original left search history to control evaluated arguments even though
the supplied step is on the right. Reversing the directed template does not
supply an unproved converse one-way property.

\paragraph{Proof strategy.}
Extract a complete local occurrence from the actual right rewrite. For
arithmetic, allocation and beta, prove its needed outer edges are retained
from the original common context. For copying, read and assignment, select
only the sharing path inspected by the rule. For equality of two references
to one store, select the union of the two inspected paths; for distinct
stores, select a path in each tree. Extend these selections to complete
occurrences of the unchanged original rule sources. Unused sibling branches
become boundary attachments rather than extra retainedness assumptions.

Reverse the filling certificates, transfer the retained occurrence, and
construct two ordinary forward machine rewrites. Prove that their computed
common residual context is binding-free, then turn its reversed-template
specimen into a specimen for the opposite of the original system. Successor
determinism identifies the initially supplied right successor with the one
constructed from the smaller occurrence. Thus rematching never substitutes
a different execution for the given step. The final exhaustive dispatch and
representation-class reflection establish the theorem.

\subsection{The opposite specimen system}
\termdef{opposite-residual-specimen}
For a directed pre-template $T$, let $T^{\mathrm{op}}(G,H)$ mean $T(H,G)$.
Output closure is symmetric: if one of the pair's graphs is one-way, the
same disjunction holds after exchanging them. Hence output closure for $T$
implies output closure for $T^{\mathrm{op}}$.

A specimen for $T^{\mathrm{op}}$ becomes an original $T$-specimen by keeping
its context and token, swapping its two filling families, and reversing
each recursive filling certificate. The closedness proofs swap with the
graphs and the allowed-context proof is unchanged. Consequently an
identity-refined quasi-specimen for $T^{\mathrm{op}}$ yields an
identity-refined quasi-specimen for the opposite of the original system.
This construction requires no symmetry of $T$.

The converse input diamond applies the generic retained-footprint diagram
to the reversed certificates. Its two steps remain ordinary machine
rewrites. The common residual graph and its binding freedom are constructed
before applying the preceding specimen conversion. The active input and
its resumed position are preserved by the complete occurrence and the
computed cut. This lemma is the common interface for recursively filled
beta and copying sources.
\lean{Copying/ConverseSpecimenTransfer}{12}{15}{PreTemplate.OutputClosed.converse}
\lean{Copying/ConverseSpecimenTransfer}{19}{31}{Specimen.unconverse}
\lean{Copying/ConverseSpecimenTransfer}{35}{41}{Specimens.converse_quasi}
\lean{Copying/ConverseSpecimenTransfer}{46}{72}{RetainedFootprint.converse_input_diamond}
\par\penalty0

\subsection{Extending selected trees to complete equality sources}
\termdef{equality-tree-rematching}
A tree occurrence used for equality must preserve each inspected input and
the root output. In the same-store case it preserves input indices $0$ and
$1$; in the distinct-store case each of the two occurrences preserves its
input $0$. These are proved properties of the selected subtrees below,
not assumptions on an arbitrary outer machine match.

For the same-store frame, map tree edges by the selected occurrence and
leave the comparator, two instance edges and store edge unchanged. The port
map fixes the operation's active input, store output and frame vertices;
it maps tree internals by the occurrence and sends each other external input
to its corresponding old tree vertex. The tree-port compatibility lemma
checks the two distinguished inputs, all remaining inputs, the root and
internal vertices separately. It then proves preservation of every source
and target row. Edge injectivity follows from the disjoint frame tags and
the tree occurrence's injectivity. Primitive labels and complete recursive
labels are preserved, yielding a full occurrence of the unchanged equality
source. Its active-input equation is definitional.
\lean{Copying/EqualityTreeRematch}{11}{20}{AtomEquality.Same.TreeRematch.ports}
\lean{Copying/EqualityTreeRematch}{22}{27}{AtomEquality.Same.TreeRematch.edges}
\lean{Copying/EqualityTreeRematch}{34}{51}{AtomEquality.Same.TreeRematch.tree_ports}
\lean{Copying/EqualityTreeRematch}{53}{89}{AtomEquality.Same.TreeRematch.embedding}
\par\penalty0
\par\penalty0
\lean{Copying/EqualityTreeRematch}{91}{94}{AtomEquality.Same.TreeRematch.occurrence}
\lean{Copying/EqualityTreeRematch}{96}{96}{AtomEquality.Same.TreeRematch.occurrence_input}
For distinct stores, use separate tree port and edge maps in the two tagged
components. The external atom inputs split into the left and right tree
blocks, so the second block is reindexed by subtracting the first block's
length. The left-tree and right-tree compatibility lemmas verify these
index calculations and the two fixed roots. The remaining frame edges
are unchanged. The same incidence, injectivity and complete-label checks
give the full two-tree occurrence and its exact active-input equation.
\lean{Copying/EqualityTreeRematch}{107}{121}{AtomEquality.Different.TreeRematch.ports}
\lean{Copying/EqualityTreeRematch}{123}{130}{AtomEquality.Different.TreeRematch.edges}
\lean{Copying/EqualityTreeRematch}{139}{154}{AtomEquality.Different.TreeRematch.left_tree_ports}
\lean{Copying/EqualityTreeRematch}{157}{173}{AtomEquality.Different.TreeRematch.right_tree_ports}
\par\penalty0
\par\penalty0
\lean{Copying/EqualityTreeRematch}{175}{223}{AtomEquality.Different.TreeRematch.embedding}
\lean{Copying/EqualityTreeRematch}{225}{228}{AtomEquality.Different.TreeRematch.occurrence}
\lean{Copying/EqualityTreeRematch}{230}{230}{AtomEquality.Different.TreeRematch.occurrence_input}
\subsection{Reversing retained local certificates without reversing computation}
The converse proof starts from an actual local rule occurrence in the
\emph{right} expanded graph $C[R]$. Reverse the recursive certificates to
obtain fillings for the directed relation $T^{\mathrm{op}}$, with the same
literal common context $C$ and the same two actual graph families in
opposite order. This does not assume that $T$ is symmetric and does not
reverse a machine transition. Instead, construct an ordinary forward
machine transition in $C[L]$ and relate the two resulting states through
the opposite of the original specimen system.

\paragraph{Generic local transfer.}
The retained local construction is generalized to an arbitrary pre-template
$T$. For beta, pad the actual paired body together with its binding-free
common context, expose its bound output, and splice the unchanged stable
argument. Template padding closure is the only additional generic
hypothesis. The stable argument's reflexive witness has its actual graph
as common context, and that graph is hole-free, so the previously proved
splice theorem establishes binding freedom of this local residual.

Restricting the decoration to the original lambda box gives the box view.
The source isomorphism retains every other original component, including
the full stable argument. In the retained case, transport the original
beta rule along this isomorphism. In the recursive case, apply the original
beta rule to the paired body and use the padded, opened and spliced
witness just constructed. Both branches produce an actual rule step with
the correct input resume coordinate.

\lean{Copying/GenericRetainedLocal}{12}{21}{RetainedLocalTransfer.betaResidual}
\lean{Copying/GenericRetainedLocal}{27}{27}{RetainedLocalTransfer.Beta.boxes}
\lean{Copying/GenericRetainedLocal}{29}{37}{RetainedLocalTransfer.Beta.sourceIso}
\lean{Copying/GenericRetainedLocal}{41}{58}{RetainedLocalTransfer.Beta.transport}
For copying, a primitive footprint is rigid, so its paired local rule is
obtained by isomorphism transport. The lambda case inspects the original
box and either retains it or duplicates its actual paired body witness.
The latter's common context is binding-free and has no shallow holes;
recursive body holes remain under boxes. Induction over every original
copyability constructor includes instance and constant payloads, active
application and arithmetic payloads, lambdas, and representation changes.
The representation case transports the source decoration and the resulting
actual rule step. This generic copying theorem needs no closure condition
on the template itself.

\lean{Copying/GenericRetainedLocal}{64}{75}{RetainedLocalTransfer.Copying.primitive_transport}
\lean{Copying/GenericRetainedLocal}{81}{81}{RetainedLocalTransfer.Copying.boxes}
\lean{Copying/GenericRetainedLocal}{83}{89}{RetainedLocalTransfer.Copying.sourceIso}
\lean{Copying/GenericRetainedLocal}{91}{110}{RetainedLocalTransfer.Copying.lambda_transport}
\par\penalty0

\par\penalty0
\lean{Copying/GenericRetainedLocal}{114}{149}{RetainedLocalTransfer.Copying.transport}
\paragraph{Actual right-occurrence extraction.}
Padding closure is preserved by reversing the pair relation: apply the
original padding theorem to the pair with its endpoints exchanged. Given
an actual complete right occurrence whose footprint avoids shallow holes,
its reversed certificates determine the retained decoration. Apply generic
copying or beta local transport, then identify the decorated source with
the exact reverse skeleton expansion using the proved host-port equation.
The output is an actual left local rule and a reversed filling witness
between the actual two reducts. Copying also supplies absence of shallow
holes in the residual context.

For beta, prove admissibility of the actual cut determined by the right
occurrence. Its reflected input port equals the original common-context
token by injectivity of the expansion host map. Reflect the original
\emph{left} rooted search to this same common-context coordinate. The
rooted-cut theorem then uses the right beta source's guarded-or-thunk
outputs and the original operational holes' term input interfaces.
Thus the proof never needs a one-way property of the right operational
filling or a new rootedness hypothesis about a synthesized source graph.

\lean{Copying/ConverseDeepExtraction}{12}{13}{PreTemplate.PaddingClosed.converse}
\lean{Copying/ConverseDeepExtraction}{24}{40}{OperationalTemplates.ConverseResidual.copying_local}
\lean{Copying/ConverseDeepExtraction}{56}{68}{OperationalTemplates.ConverseResidual.beta_local}
\lean{Copying/ConverseDeepExtraction}{75}{101}{OperationalTemplates.ConverseResidual.beta_frame_binding}
\subsection{A retained minimal spine for every actual copying occurrence}
The supplied copying tree may contain irrelevant sibling branches whose
edges originate in shallow right fillings. The proof rematches the actual
source with its selected contraction spine and its complete payload.
A counted adjacency trace gives an uncounted ancestry path by induction
on the trace. The already proved minimal occurrence has an exact selected
input, fixes the original payload coordinates, and gives a counted path
from that input's contraction to every retained spine edge.

Map these paths through the supplied whole-graph occurrence. The active
contraction is retained because the common-context token is not entering
a hole. Binding freedom then excludes shallow holes along each selected
spine path. Every payload edge is likewise excluded by the proved
copying-payload barrier, including the lambda's boxed tail. Together these
arguments establish avoidance for the complete smaller footprint; they
make no assertion about uninspected sibling branches.

For an arbitrary copyable payload, first use its proved recursive
representation isomorphism to a normal form, construct the smaller
occurrence there, and compose with the supplied original occurrence.
The result is an actual full-label occurrence of an original copying
source at the exact same host token, with proved copyability and proved
footprint avoidance. Normalization changes only the proof's representation
of the matched source, not the actual machine or its transition rules.

\lean{Copying/ConverseCopyFootprint}{17}{21}{BindingOverlap.trace_ancestor}
\lean{Copying/ConverseCopyFootprint}{25}{77}{BindingOverlap.normal_minimal_copy_occurrence}
\lean{Copying/ConverseCopyFootprint}{82}{94}{BindingOverlap.minimal_copy_occurrence}
\paragraph{Universal converse copying case.}
Apply this minimal-spine theorem to the reversed certificates. Extract the
actual left local copy step and the reversed residual witness. The original
copyable source guards every output attachment; the residual is binding-free
and has no shallow holes, so the open guarded graft theorem proves binding
freedom of the actual common-context replacement. The converse input diamond
constructs both whole-graph rewrites and an equality-up-to quasi-specimen
for the \emph{opposite of the original system}.

The quantitative theorem accepts any independently supplied right
successor. Proved successor determinism identifies it with the constructed
right reduct. It uses zero additional right steps and one actual left
rewrite, then the original opposite quasi-specimen. The theorem requires
neither a left match nor a retained footprint certificate from the caller;
it covers every original finite tree presentation and copyable payload.

\lean{Copying/ConverseCopyResidual}{19}{37}{OperationalTemplates.ConverseResidual.copying_diamond}
\lean{Copying/ConverseCopyResidual}{41}{55}{OperationalTemplates.ConverseResidual.copying}
\subsection{Converse arithmetic, allocation and arbitrary deep beta}
Unary and binary arithmetic use the original left search history to prove
that the actual right head and its evaluated numeric operands all belong
to retained common-context edges. The generic sink diamond applies to the
reversed certificates. Allocation uses the same readiness argument for its
unchanged stable operand, then the guarded-or-thunk replacement theorem
preserves its full output interface. Stable lambdas and instances remain
included. In each case, convert the residual for reversed certificates
into an opposite-system residual for the original directed template;
then apply the proved matching quantitative argument to the independently
supplied right successor.

\lean{Copying/ConverseRetainedResiduals}{20}{36}{OperationalTemplates.ConverseResidual.unary}
\lean{Copying/ConverseRetainedResiduals}{39}{56}{OperationalTemplates.ConverseResidual.binary}
\lean{Copying/ConverseRetainedResiduals}{60}{78}{OperationalTemplates.ConverseResidual.allocation}
For beta, the complete right footprint is derived from the original left
search history, including both application arguments and the actual
lambda box. The reversed local extraction allows an arbitrary recursively
filled body and retains the actual stable argument. The rooted cut proof
above establishes admissibility; plugging the actual local body residual
into that cut then yields the original opposite quasi-specimen at the
exact resume ports. The quantitative corollary again takes zero additional
right steps and one left rewrite. These public theorems require actual
right local occurrences, which every original contextual rewrite supplies;
no left occurrence, overlap assumption, framed decomposition or residual
admissibility premise remains.

\lean{Copying/ConverseRetainedResiduals}{92}{107}{OperationalTemplates.ConverseResidual.beta_diamond}
\lean{Copying/ConverseRetainedResiduals}{111}{125}{OperationalTemplates.ConverseResidual.beta}
\subsection{Retained atom-equality transfer}
The final transfer for each atom-equality rule starts with an actual retained
right occurrence, whose footprint is supplied by the separate rematching
proof. The same-store rule retains the true result and its one store output;
the distinct-store rule retains the false result and both distinct store
outputs. Apply the generic guarded diamond to the reversed recursive
certificates, using the original source's proved output guards and the
actual reduct's focus and hole freedom. Convert its quasi-specimen back to
the opposite of the original system. Proved determinism then treats any
independently supplied right successor. These two helpers require footprint
avoidance; the complete equality theorems discharge that premise using the
selected-subtree construction.

\lean{Copying/ConverseEqualityTransfer}{17}{35}{OperationalTemplates.ConverseResidual.retained_same}
\lean{Copying/ConverseEqualityTransfer}{39}{57}{OperationalTemplates.ConverseResidual.retained_different}
\subsection{Extracting the two selected atom ancestries}
\termdef{selected-contraction-subtree}
The converse equality argument must retain the paths selected by the two
already evaluated reference instances. An arbitrary presentation of the original
atom tree may also include branches which lie inside a shallow template filling.
Those branches cannot simply be declared to belong to the common context. We
instead construct a smaller tree containing the selected paths and expose every
other attachment as a boundary input. This is a local occurrence of an unchanged
machine rule: its selected atom inputs and its store output have the original
coordinates, and its sharing edges retain their ordered incidences.

\paragraph{Forward-closed selections.}
Let $T$ be a finite contraction tree of sharing sort $s$. A selection consists
of a predicate $P$ on its edges, closed under forward adjacency, together with
$P(r)$ for the root edge $r$. The selected vertices are exactly those incident
to a selected edge; the selected edges are the edges satisfying $P$.
The vertex and edge carriers therefore remain finite subsets of the original
carriers. The ordered source and target lists are inherited without changing
the underlying ports. Lean packages these data as
\lean{Copying/SelectedContractionSubtree}{10}{13}{ContractionTree.Selection}
\lean{Copying/SelectedContractionSubtree}{18}{19}{Selection.Incident}
\lean{Copying/SelectedContractionSubtree}{21}{21}{Selection.Vertex}
\lean{Copying/SelectedContractionSubtree}{22}{22}{Selection.Edge}
\par\penalty0
The incidence definitions are
\lean{Copying/SelectedContractionSubtree}{24}{25}{Selection.source}
\lean{Copying/SelectedContractionSubtree}{26}{27}{Selection.target}
Their elements carry proofs of membership in the selected vertex set. Erasing
these proofs gives exactly the original ordered source and target lists:
\lean{Copying/SelectedContractionSubtree}{29}{30}{Selection.source_val}
\lean{Copying/SelectedContractionSubtree}{31}{32}{Selection.target_val}
\par\penalty0
Consequently, membership in either selected incidence list is equivalent to
membership in the corresponding original list. The proofs map a membership
witness through the subtype projection in one direction and restore its
incident-vertex certificate in the other:
\lean{Copying/SelectedContractionSubtree}{34}{41}{Selection.mem_source}
\lean{Copying/SelectedContractionSubtree}{42}{49}{Selection.mem_target}
A selected vertex is a new input precisely when no selected edge produces it.
The output remains the original sole output; the selected root certifies its
presence. These definitions are
\lean{Copying/SelectedContractionSubtree}{51}{51}{Selection.Input}
\lean{Copying/SelectedContractionSubtree}{53}{53}{Selection.output}
Every selected vertex has wire type $s$, since it occurs in an original
sharing-edge source or target. In the source case a weakening is impossible,
while a contraction has only $s$-typed inputs; both possible labels have one
$s$-typed output. This proves
\lean{Copying/SelectedContractionSubtree}{55}{67}{Selection.label}
\par\penalty0

Every selected vertex other than the original output is consumed by a selected
edge. If it already occurs in a selected source there is nothing to prove.
Otherwise it is produced by a selected edge, and original source coverage
provides its consumer. Forward closure selects that consumer. Conversely, the
original graph prohibits consuming the output. Thus the output is exactly the
only vertex lacking a source incidence. The root produces this output, so it
cannot also be an input. These boundary facts are
\lean{Copying/SelectedContractionSubtree}{69}{87}{Selection.source_coverage}
\lean{Copying/SelectedContractionSubtree}{89}{91}{Selection.output_not_input}
Choose a finite list containing each new input exactly once. The resulting
set presentation has that input order, the original output, the selected
carriers, and the inherited incidence lists. Distinctness and uniqueness of
incidence follow from their original counterparts; the two boundary-coverage
laws follow from the preceding argument and the definition of a new input.
All typing obligations follow by erasing the subtype certificates. This is the
fully verified graph construction
\lean{Copying/SelectedContractionSubtree}{94}{162}{Selection.setNet}
It makes a classical finite choice of presentation, without adding a transition
or changing any operation.

\paragraph{The actual local occurrence.}
Normalizing the set presentation gives the usual tagged internal/input/output
ports. Decode such a port and then forget its membership certificate to obtain
its original vertex. Applying this map to either normalized incidence row
recovers the original row, because encoding followed by decoding is the
identity. An adjacency witness therefore also maps to an original adjacency
witness. The relevant declarations are
\lean{Copying/SelectedContractionSubtree}{168}{169}{Selection.ports}
\lean{Copying/SelectedContractionSubtree}{171}{176}{Selection.source_ports}
\lean{Copying/SelectedContractionSubtree}{178}{183}{Selection.target_ports}
\lean{Copying/SelectedContractionSubtree}{185}{193}{Selection.adjacent}
\par\penalty0
The normalized graph is itself a contraction tree. Its labels are inherited;
a cyclic selected path would map injectively on edges to a cyclic original
path, contradicting the original acyclicity proof. The inclusion preserves
all edge labels and ordered incidences, and hence gives a complete local
occurrence at every box depth. It sends input position $i$ to the $i$th chosen
boundary vertex and sends output zero to the original output zero:
\lean{Copying/SelectedContractionSubtree}{196}{205}{Selection.tree}
\lean{Copying/SelectedContractionSubtree}{208}{217}{Selection.occurrence}
\lean{Copying/SelectedContractionSubtree}{219}{221}{Selection.occurrence_input}
\lean{Copying/SelectedContractionSubtree}{223}{224}{Selection.occurrence_output}
\par\penalty0

\paragraph{Two prescribed inputs.}
A finite selected boundary can be ordered with any two distinct prescribed
inputs $a,b$ first. Enumerate all selected inputs, remove repetitions, filter
out $a$ and $b$, and prepend them in that order. The remaining list has no
repetitions, and the final list contains exactly the inputs. The local
finite-list induction removes an element only when it already appears in the
inductively obtained list, thereby preserving membership. This proves
\lean{Copying/TwoInputContractionSubtree}{31}{62}{Selection.two_input_order}
the elementary finite-list proof is
\lean{Copying/TwoInputContractionSubtree}{9}{25}{distinct_members}
For an original input $i$, an edge is \emph{from $i$} if the input is consumed
by some edge $a$ and a finite counted adjacency trace runs from $a$ to that
edge. Select the union of the predicates for two inputs $i,j$. Appending one
adjacency to a trace proves forward closure; the original input-to-root theorem
shows that the root is selected. Each specified input remains a new boundary
input: its original consumer is selected, and no original edge can target an
original input. These facts are formalized by
\lean{Copying/TwoInputContractionSubtree}{68}{69}{ContractionTree.FromInput}
\lean{Copying/TwoInputContractionSubtree}{78}{85}{ContractionTree.twoSelection}
\lean{Copying/TwoInputContractionSubtree}{88}{99}{ContractionTree.twoSelection_input}
\par\penalty0
The elementary trace-extension induction is
\lean{Copying/TwoInputContractionSubtree}{71}{75}{trace_last}
\begin{theorem}[Two-input subtree extraction]
For distinct input coordinates $i,j$ of any finite contraction tree $T$ and
any depth $d$, there are $h$, a contraction tree $U$ with $h+2$ inputs, and a
complete local occurrence $f:U\hookrightarrow T$ at depth $d$ such that
$f(0)=i$, $f(1)=j$, and the sole output is fixed. Every selected edge $f(e)$
has a counted original ancestry from $i$ or from $j$.
\end{theorem}
Take the preceding union selection, order its new inputs with $i,j$ first,
and apply the normalized-tree occurrence construction. The ancestry assertion
is the stored selection predicate of each selected edge. No restriction on
unselected siblings is assumed:
\lean{Copying/TwoInputContractionSubtree}{103}{134}{ContractionTree.two_input_subtree}
\paragraph{Same-store equality on the converse side.}
Consider an original binding-free operational specimen with expansions
$C[L]$ and $C[R]$, at a non-entering rewrite token. Assume the original left
expansion is rooted. Suppose an actual equality-rule occurrence in $C[R]$
uses the same store for both operands. Rooted argument readiness retains both
selected instances in the common context. Each is a binding edge, so binding
freedom retains every edge downstream of either instance.

Apply two-input extraction to the occurrence's original atom tree and extend
that occurrence through the unchanged equality frame. Its comparison edge
is retained because the token is non-entering. Its two instances are retained
by the preceding argument. Each selected sharing edge is downstream of one
of those instances, by its stored original ancestry. Finally, the common store
is downstream of the left instance by the original successful resolution
trace. This proves that the \emph{entire rematched footprint} avoids shallow
holes, while making no such assertion about unused branches in the original
larger rule presentation:
\lean{Copying/ConverseSameEquality}{20}{81}{ConverseOverlap.same_subtree_avoids}
The extraction theorem then discharges the selected-subtree premise for every
actual original same-store occurrence and preserves the original comparison
token. Thus it constructs a complete retained occurrence of an original
same-store rule, rather than assuming the desired left match:
\lean{Copying/ConverseSameEquality}{96}{106}{ConverseResidual.same_occurrence}
The existing guarded primitive residual construction now supplies actual
right and left rewrites and an admissible opposite quasi-specimen. Successor
coherence identifies the constructed right result with any independently
supplied actual successor. The quantitative matching uses zero further
right steps and one left step, so the cost comparison is $Q(1+0,1)$ by
reflexivity. Neither a residual-admissibility certificate nor an opposite-step
existence assumption remains:
\lean{Copying/ConverseSameEquality}{112}{123}{ConverseResidual.same}
\Needspace{10\baselineskip}
\subsection{Rematching the selected read and assignment trees}
A converse read or assignment can be presented with an atom contraction tree
containing irrelevant sibling branches. A branch of that supplied tree may
belong to a right-hand shallow template filling. Consequently, the whole
supplied tree cannot simply be declared retained. The proof instead extracts
its actual selected input-to-store spine, leaving unused branches attached at
new boundary inputs. The resulting smaller graph is still a contraction tree
and is still an instance of the original local rule. No computation rule,
stability condition, term or translation changes.

First consider arbitrary atom trees $U$ and $T$ and a complete occurrence
$f:U\hookrightarrow T$. Assume that $f$ sends the distinguished first input
to the first input and the root output to the root output. Other inputs may
attach to internal wires of $T$. The read rematching construction extends $f$
to an occurrence of the entire read source. Its port function preserves the
active input, store-content output and the three frame wires; it maps an
internal tree wire through $f$, and maps every additional input through the
corresponding tree attachment. Its edge function preserves the reader,
instance and store edges and uses $f$ on tree edges. The tree-port lemma
proves that insertion commutes with this map. Internal wires use the
occurrence's internal-coordinate equation; the first input and root output
use the two assumptions; every other input uses its explicit ordered frame
attachment.

\lean{Copying/ReadTreeRematch}{11}{19}{Dereference.TreeRematch.ports}
\lean{Copying/ReadTreeRematch}{21}{25}{Dereference.TreeRematch.edges}
\lean{Copying/ReadTreeRematch}{31}{47}{Dereference.TreeRematch.tree_ports}
The incidence-preserving embedding follows by cases on the four kinds of
read-source edge. The three fixed primitive edges have unchanged rows.
For a tree edge, map composition and the tree-port lemma transform the
source and target rows using the original occurrence equations. Edge
injectivity follows from the disjoint component tags and injectivity of $f$.
The complete occurrence constructor then supplies the induced injective
internal map and full recursive-label preservation; it does not require
injectivity of all attached boundary ports. Finally, its active input is
literally the original read input.

\lean{Copying/ReadTreeRematch}{49}{83}{Dereference.TreeRematch.embedding}
\lean{Copying/ReadTreeRematch}{85}{88}{Dereference.TreeRematch.occurrence}
\lean{Copying/ReadTreeRematch}{90}{90}{Dereference.TreeRematch.occurrence_input}
Assignment uses the same construction with four frame wires and an
additional value component. The stable new-value graph $G$ is kept literally:
all of its edges and internal wires are fixed, and each of its ordered outputs
keeps the same external position. The old-content output is fixed separately.
Thus rematching neither replaces the new value by an observationally related
value nor changes which store receives it. The port and edge maps extend the
tree occurrence and preserve this entire value component. Their tree-port
lemma has the same input, internal-wire and output proof as for reading.

\lean{Copying/AssignmentTreeRematch}{11}{22}{Assignment.TreeRematch.ports}
\lean{Copying/AssignmentTreeRematch}{24}{31}{Assignment.TreeRematch.edges}
\lean{Copying/AssignmentTreeRematch}{37}{59}{Assignment.TreeRematch.tree_ports}
For the assignment embedding, the only additional incidence case is an
edge of $G$. Its source and target lists commute with insertion because the
value's sole input attaches to the fixed new-value frame wire, its internal
wires are fixed, and its outputs attach in their original order. The complete
occurrence again follows from full label and incidence preservation, and its
active input is unchanged. These structural constructors work for arbitrary
$G$; the actual assignment rule later retains its original proof that $G$ is
stable.

\lean{Copying/AssignmentTreeRematch}{61}{136}{Assignment.TreeRematch.embedding}
\lean{Copying/AssignmentTreeRematch}{138}{143}{Assignment.TreeRematch.occurrence}
\lean{Copying/AssignmentTreeRematch}{145}{145}{Assignment.TreeRematch.occurrence_input}
\Needspace{10\baselineskip}
\subsection{The original binding barrier retains each selected store spine}
Let $C$ be an original binding-free operational specimen context with paired
fillings $L,R$. Suppose the common rewrite token is not entering a shallow
hole, and the left endpoint has a rooted search history to that token.
Consider an active operation in $C[R]$ and a binding edge $a$ immediately
following one of its arguments. The argument-readiness theorem first proves
that the original owner of $a$ is not a shallow hole. This yields its retained
corolla certificate. Binding-label reflection supplies the binding property
of the original edge in $C$, and reflection of the argument adjacency supplies
a source witness for it.

If $b$ is any actual descendant of $a$, the original binding-free barrier
reflects that ancestry into $C$. Were the original owner of $b$ a hole, it
would differ from the already non-hole owner of $a$; the reflected ancestry
would therefore be positive. This contradicts the precise binding-free
condition, which prohibits a path from a sourced binding edge to a hole.
The proof applies to each selected spine independently and never assumes
that every sibling in the supplied tree is retained.

\lean{Copying/ConverseBindingDescendant}{12}{35}{OperationalTemplates.ConverseOverlap.argument_binding_descendant_not_hole}
For a read subtree, require that each of its mapped edges is reached by
an actual counted path from a contraction consuming the original first atom
input. The original instance enters that contraction, so each such edge is
a descendant of the instance. The entire selected source avoids shallow
holes: the reader is retained because the token is non-entering; the instance
uses reflexive ancestry in the binding-argument lemma; tree edges use the
selected paths; and the store uses the original successful store-resolution
trace. This establishes avoidance for the complete rematched occurrence,
including its operation and store edges.

The assignment argument is identical for the destination instance, selected
tree and store. Every edge of the original stable value consumes its sole term
input, so it is adjacent to the active assignment head at the new-value
argument. Rooted argument readiness excludes a shallow-hole origin for these
edges as well. This preserves all stable operand forms and their full output
interfaces, including thunk outputs.

\lean{Copying/ConverseStoreFootprint}{19}{53}{OperationalTemplates.ConverseOverlap.read_subtree_avoids}
\lean{Copying/ConverseStoreFootprint}{57}{99}{OperationalTemplates.ConverseOverlap.assignment_subtree_avoids}
\Needspace{10\baselineskip}
\subsection{Actual converse read and assignment residuals}
The public read-occurrence theorem starts with any actual complete right read
occurrence. The minimal-spine theorem extracts a finite oriented spine from
the supplied atom tree, fixes its selected first input and root output, and
provides the counted ancestry of every retained spine edge. Extend that tree
occurrence through the read frame and compose it with the actual host
occurrence. The preceding avoidance theorem discharges the entire footprint
condition. The public assignment-occurrence theorem performs the same steps
while preserving the literal stable value graph. Their conclusions contain
both an actual complete source occurrence and its exact original rewrite-token
attachment. No smaller-tree certificate is assumed by either caller.

\lean{Copying/ConverseStoreResiduals}{22}{35}{OperationalTemplates.ConverseResidual.read_occurrence}
\lean{Copying/ConverseStoreResiduals}{39}{55}{OperationalTemplates.ConverseResidual.assignment_occurrence}
The two diamond theorems now apply the already proved primitive transfer
construction with the paired filling certificates reversed. For reading,
the original sourced-instance ancestry guards the content output, and the
literal read reduct is focus-free and hole-free. For assignment, the original
source has binding-or-thunk outputs. Hole typing follows from the closed
operational filling, and stability proves the required safety of the literal
assignment reduct. The guarded or thunk-aware replacement theorem establishes
binding freedom of the resulting common residual context; it is not an extra
hypothesis of either diamond.

Each diamond constructs a genuine right rewrite and a genuine left rewrite
and gives an equality-up-to quasi-specimen between their resumed search
states. The reversed-template quasi-specimen is then converted into a
quasi-specimen for the opposite of the \emph{original} directed specimen
system. This conversion requires no symmetry of the template. All original
unmatched shallow and deep fillings remain in the common residual.

\lean{Copying/ConverseStoreResiduals}{61}{74}{OperationalTemplates.ConverseResidual.read_diamond}
\lean{Copying/ConverseStoreResiduals}{78}{93}{OperationalTemplates.ConverseResidual.assignment_diamond}
Finally, the quantitative theorems quantify over any supplied actual right
machine successor. Successor determinism identifies it with the constructed
right successor in the representation quotient. There are zero further steps
on that side and one original left rewrite on the other, so the cost comparison
is $Q(1+0,1)$ and follows from reflexivity of an arbitrary cost order $Q$.
The equality-up-to residual implies the required refinement-up-to residual
for arbitrary orders $P$ and $U$. Neither theorem assumes a left rule match,
a shallow-overlap classification, a residual-admissibility certificate or a
canonical shape for the supplied right successor. The original left rootedness
and the non-entering common token are the operational specimen hypotheses
used to construct these facts.

\lean{Copying/ConverseStoreResiduals}{96}{108}{OperationalTemplates.ConverseResidual.read}
\lean{Copying/ConverseStoreResiduals}{111}{124}{OperationalTemplates.ConverseResidual.assignment}
\Needspace{10\baselineskip}
\subsection{Distinct-store equality uses two independent selected spines}
The distinct-store equality rule has two instances, two atom trees and two
stores. Its rematched source contains only the selected spine in each tree,
while preserving both actual store edges and their ordered content outputs.
For each instance separately, the rooted argument theorem establishes its
original non-hole owner, and the binding-descendant barrier excludes shallow
holes along its selected spine. The corresponding successful resolution trace
retains its store. The comparison head is retained at the non-entering token.
A case split on the seven component tags therefore proves avoidance for the
complete rematched distinct-store source. In particular, the proof does not
conflate store identity with the shape of either tree.

\lean{Copying/ConverseDifferentEquality}{19}{93}{OperationalTemplates.ConverseOverlap.different_subtrees_avoids}
The occurrence theorem invokes minimal-spine extraction independently for
the two original trees, extends both complete occurrences through the literal
distinct-store comparison frame, and composes with the original right match.
Both stores remain distinct because complete occurrence maps are injective
on edge occurrences. The exact comparison token and both content outputs
remain attached as before. Thus the newly constructed source is an actual
instance of the original false-result rule, rather than a new decision rule
inferred solely from an intrinsic selector.

\lean{Copying/ConverseDifferentEquality}{108}{124}{OperationalTemplates.ConverseResidual.different_occurrence}
The final distinct-store theorem applies the retained equality transfer to
this constructed occurrence. That transfer builds both original computations,
uses the source's binding ancestors to establish residual admissibility, and
returns an opposite quasi-specimen. The arbitrary supplied right successor
is handled by the same successor-determinism argument as reading and
assignment. Its quantitative comparison is again the exact one-step match
$Q(1,1)$ for every cost order. The same-store equality branch needs a union of
two selected paths in one tree and is proved separately; independent spine
extraction is used here only for the genuinely distinct-store rule.

\lean{Copying/ConverseDifferentEquality}{130}{141}{OperationalTemplates.ConverseResidual.different}
\subsection{Exhaustive assembly of universal converse robustness}
The local-case theorem performs case analysis on the original local-rule
constructor. Allocation, both arithmetic operations, read, assignment and
both equality decisions use their now unconditional converse residual
lemmas. Beta uses the original left rootedness to prove its complete cut
admissible; copying uses the selected spine and retained recursive payload.
Each branch supplies the same quantitative conclusion for the independently
provided right successor. The original stability and copyability proofs
are passed unchanged.

An arbitrary contextual rewrite supplies its own complete local occurrence,
including its literal source and target rule and its exact active token.
Applying the local-case theorem therefore discharges all match extraction
premises. At the representation-class level, reflect a successor to an actual
finite step. A search step is impossible in rewrite mode, leaving exactly the
rewrite case already established.

Finally, take an arbitrary eligible specimen in the opposite of the original
system. Its original left state supplies the rooted search history. Rewrite
its endpoints into the common-context host coordinates, transport the
non-entering premise to that token, and apply the preceding class-step
result. The resulting quantitative matching alternative is precisely the
original opposite-system robustness conclusion. Thus the proof is universal
over specimens, contexts, depths and right successors, with no residual
existence assumption hidden in the final theorem.
\lean{Copying/FullOperationalConverseRobustness}{23}{43}{OperationalTemplates.ConverseResidual.local_case}
\lean{Copying/FullOperationalConverseRobustness}{47}{58}{OperationalTemplates.ConverseResidual.rewrite_case}
\lean{Copying/FullOperationalConverseRobustness}{62}{75}{OperationalTemplates.ConverseResidual.class_step}
\lean{Copying/FullOperationalConverseRobustness}{80}{109}{OperationalTemplates.template_robust_converse}
\subsection{Concrete contextual consequences for both operational families}
\termdef{operational-contextual-equivalence}
The complete forward and converse proofs now discharge the robustness
hypotheses for micro-beta and reference creation. The existing input-safety
proof consumes their finite search-and-compute prefixes, preserving all
stability premises. Therefore neither safety nor robustness is an assumption
of the following semantic conclusions.

First, a recursively filled closed pair whose actual common context is
binding-free supplies a concrete initial lifting. Construct the original
specimen at the common input in search mode, use the actual endpoint
isomorphisms to identify its two states, and establish rootedness by the
zero-length initial search. This construction is generic in the pre-template.

\lean{Copying/OperationalContextualEquivalence}{14}{29}{FillingWitness.binding_initial_lifting}
For any operational pair of open programs with the same variable and atom
interface, apply the concrete contextual robustness theorem. In the forward
direction, use the improvement cost order: every successful observation of
the left graph has a successful observation of the right graph with no
larger transition count. In both directions, use the universal cost relation
to obtain contextual termination equivalence. Its equivalent formulation
states that the two actual closed programs terminate together in every
original binding-free closing context. The context's surrounding depth is
unrestricted, and neither a terminating continuation nor any syntactic
restriction on the surrounding language is assumed.

\lean{Copying/OperationalContextualEquivalence}{35}{40}{OperationalTemplates.pair_contextual_improvement}
\lean{Copying/OperationalContextualEquivalence}{44}{50}{OperationalTemplates.pair_contextual_equiv}
\lean{Copying/OperationalContextualEquivalence}{54}{57}{OperationalTemplates.pair_contextual_termination}
The same argument applies directly to arbitrary recursive filling witnesses,
including many simultaneous and deeply boxed operational occurrences.
Apply the abstract sufficiency theorem to the initial lifting and the proved
forward robustness for improvement. For equivalence, apply forward and
converse sufficiency with the universal cost relation. Unpacking the two
refinements yields equivalence of existence of successful finite runs.
The witnesses' transition counts may differ: the operational search and
compute prefixes have not been made cost-free. These results establish the
micro-beta and reference-creation steps needed by Figure~22; they do not
assert the full Proposition~12.2 chain or its remaining templates.

\lean{Copying/OperationalContextualEquivalence}{61}{68}{OperationalTemplates.filling_improvement}
\lean{Copying/OperationalContextualEquivalence}{72}{83}{OperationalTemplates.filling_equiv}
\lean{Copying/OperationalContextualEquivalence}{87}{99}{OperationalTemplates.filling_termination}